\documentclass[11pt]{article}
\usepackage[a4paper, left=3cm, top=3cm, right=2cm, bottom=2cm]{geometry}
\usepackage{setspace}
\usepackage{amsmath, amssymb, amsthm}
\usepackage{fontspec}
\setmainfont{Arial}
\singlespacing

\begin{document}
    Redes complexas são fundamentais para modelar sistemas interconectados em diversas áreas da ciência e da computação, permitindo descrever desde processos de difusão e transporte de partículas até a propagação de epidemias e algoritmos de busca. Compreender como processos se comportam assintoticamente é crucial para prever o estado de equilíbrio do sistema. Enquanto simular passeios aleatórios em redes de larga-escala apresenta um alto custo computacional, a teoria espectral de grafos oferece uma abordagem analítica rigorosa, conectando a topologia combinatória do grafo às propriedades espectrais de suas matrizes associadas.
    Neste trabalho, apresentamos uma revisão teórica que explora a interseção entre a teoria espectral e passeios aleatórios. Especificamente, estudamos como a matriz de um passeio aleatório \(\mathbf{W}\) pode ser simetrizada pela matriz de adjacência normalizada \(\mathbf{A}\) ou pela matriz laplaciana normalizada \(\mathbf{N}\) para avaliar o comportamento estacionário de um passeio sem realizar simulações numéricas extensivas.
    Seguindo o tratamento da literatura contemporânea sobre teoria espectral e algébrica de grafos, analisamos o vetor da distribuição de probabilidade \(\mathbf{p}_t\) e descrevemos a formalização do passeio aleatório como uma iteração de potências da matriz estocástica, \(\mathbf{p}_t = \mathbf{W}^t \mathbf{p}_0\). Utilizamos transformações de similaridade para relacionar a matriz não simétrica \(\mathbf{W}\) com a matriz simétrica \(\mathbf{A}\), possibilitando o uso do Teorema Espectral para caracterizar o comportamento assintótico do passeio aleatório.
    A revisão teórica mostra que a distribuição estacionária de um passeio aleatório é governada explicitamente pelo autovetor dominante correspondente ao autovalor principal \(\lambda_1 = 1\) de \(\mathbf{A}\). Além disso, a taxa de convergência (ou tempo de mistura) do passeio aleatório é limitada pelo segundo maior autovalor \(\lambda_2\). Isso demonstra como simulações em sistemas complexos podem ser previstas por propriedades algébricas do espectro de um grafo, provendo uma fundamentação teórica para ferramentas computacionais como o PageRank do Google.
\end{document}